\documentclass[12pt, a4paper]{article}
\usepackage[UTF8]{ctex}
\usepackage{amsmath, amssymb, amsthm}
\usepackage{geometry}

\geometry{left=2.5cm, right=2.5cm, top=2.5cm, bottom=2.5cm}

\newcommand{\RR}{\mathbb{R}}
\newcommand{\e}{\mathrm{e}}

\begin{document}
\section*{11.2}
\textbf{1.}确定下列函数的自然定义域：

(3)\(u = \sqrt{R^2 - x^2 - y^2 - z^2} + \sqrt{x^2 + y^2 + z^2 - r^2} \quad (R > r)\).

解：
\[
\begin{cases}
R^2 - x^2 - y^2 - z^2 \geq 0 \\
x^2 + y^2 + z^2 - r^2 \geq 0
\end{cases}, \Rightarrow r^2 \leq x^2 + y^2 + z^2 \leq R^2.
\]

(4)\(u = \arcsin \dfrac{z}{x^2 + y^2}\).

解：
\[x^2 + y^2 \neq 0, -1 \leq \dfrac{z}{x^2 + y^2} \leq 1 \Rightarrow x^2 + y^2 \neq 0, - (x^2 + y^2) \leq z \leq x^2 + y^2.\]

\textbf{3.}若函数
\[z(x, y) = \sqrt{y} + f(\sqrt{x} - 1),\]
且当\(y = 4\)时\(z = x + 1\), 求\(f(x)\)和\(z(x, y)\).

解：由\(y = 4\)时\(z = x + 1\)，得
\[f(\sqrt{x} - 1) = x - 1.\]
令\(\sqrt{x} - 1 = t\)，则\(x = (t + 1)^2 = t^2 + 2t + 1\)，
\[f(t) = t^2 + 2t, \quad t \geq -1.\]
所以
\[z(x, y) = \sqrt{y} + f(\sqrt{x} - 1) = \sqrt{y} + (\sqrt{x} - 1)^2 + 2(\sqrt{x} - 1) = x + y - 1.\]

\textbf{5.}对多元函数证明其极限唯一性、局部有界性、局部保序性和局部夹逼性.

解：

(1)极限唯一性：设\(\lim\limits_{x \to x_0} f(x) = A\)，\(\lim\limits_{x \to x_0} f(x) = B\)，则\(\forall \varepsilon > 0\)，\(\exists \delta_1 > 0\)，当\(0 < \|x - x_0\| < \delta_1\)时，有\(|f(x) - A| < \varepsilon\)，\(\exists \delta_2 > 0\)，当\(0 < \|x - x_0\| < \delta_2\)时，有\(|f(x) - B| < \varepsilon\)，令\(\delta = \min\{\delta_1, \delta_2\}\)，则当\(0 < \|x - x_0\| < \delta\)时，有
\[|A - B| = |A - f(x) + f(x) - B| \leq |A - f(x)| + |f(x) - B| < 2\varepsilon.\]
所以\(A = B\). 

(2)局部有界性：设\(\lim\limits_{x \to x_0} f(x) = A\)，则\(\forall \varepsilon > 0\)，\(\exists \delta > 0\)，当\(0 < \|x - x_0\| < \delta\)时，有\(|f(x) - A| < \varepsilon\)，即
\[|f(x)| < |A| + \varepsilon.\]

(3)局部保序性：设\(\lim\limits_{x \to x_0} f(x) = A\)，\(\lim\limits_{x \to x_0} g(x) = B\)，且在\(O(x_0, \delta_0)\)内，\(f(x) \leq g(x)\)，则\(\forall \varepsilon > 0\)，\(\exists \delta_1 > 0\)，当\(0 < \|x - x_0\| < \delta_1\)时，有\(|f(x) - A| < \varepsilon\)，\(\exists \delta_2 > 0\)，当\(0 < \|x - x_0\| < \delta_2\)时，有\(|g(x) - B| < \varepsilon\)，令\(\delta = \min\{\delta_0, \delta_1, \delta_2\}\)，则当\(0 < \|x - x_0\| < \delta\)时，有
\[A - \varepsilon < f(x) \leq g(x) < B + \varepsilon.\]
因此\(A \leq B\).

(4)局部夹逼性：设\(\lim\limits_{x \to x_0} f(x) = \lim\limits_{x \to x_0} g(x) = A\)，且在\(O(x_0, \delta_0)\)内，\(f(x) \leq h(x) \leq g(x)\)，则\(\forall \varepsilon > 0\)，\(\exists \delta_1 > 0\)，当\(0 < \|x - x_0\| < \delta_1\)时，有\(|f(x) - A| < \varepsilon\)，\(\exists \delta_2 > 0\)，当\(0 < \|x - x_0\| < \delta_2\)时，有\(|g(x) - A| < \varepsilon\)，令\(\delta = \min\{\delta_0, \delta_1, \delta_2\}\)，则当\(0 < \|x - x_0\| < \delta\)时，有
\[A - \varepsilon < f(x) \leq h(x) \leq g(x) < A + \varepsilon.\]
故\(\lim\limits_{x \to x_0} h(x) = A\).

\textbf{6.}对多元函数证明极限的四则运算法则：假设当\(x\)趋于\(x_0\)时函数\(f(x)\)和\(g(x)\)的极限存在，则

(3)\(\lim\limits_{x \to x_0} (f(x) / g(x)) = \lim\limits_{x \to x_0} f(x) / \lim\limits_{x \to x_0} g(x) \quad (\lim\limits_{x \to x_0} g(x) \neq 0)\).

解：设\(\lim\limits_{x \to x_0} f(x) = A\)，\(\lim\limits_{x \to x_0} g(x) = B\)，且\(B \neq 0\)，则\(\forall \varepsilon > 0\)，\(\exists \delta_1 > 0\)，当\(0 < \|x - x_0\| < \delta_1\)时，有\(|f(x) - A| < \varepsilon\)，\(\exists \delta_2 > 0\)，当\(0 < \|x - x_0\| < \delta_2\)时，有\(|g(x) - B| < \varepsilon\)，令\(\delta = \min\{\delta_1, \delta_2\}\)，则当\(0 < \|x - x_0\| < \delta\)时，有
\[\lim_{x \to x_0} \dfrac{f(x)}{g(x)} - \frac{A}{B} = \frac{B(f(x) - A) - A(g(x) - B)}{Bg(x)}.\]
\[\left|\lim_{x \to x_0} \dfrac{f(x)}{g(x)} - \frac{A}{B}\right| = \left|\frac{B(f(x) - A) - A(g(x) - B)}{Bg(x)}\right| \leq \frac{|B|\varepsilon + |A|\varepsilon}{|B||g(x)|}.\]
所以\(\lim\limits_{x \to x_0} (f(x) / g(x)) = A / B\).

\textbf{7.}求下列极限：

(2)
\begin{align*}
\lim_{(x, y) \to (0, 0)} \frac{1 + x^2 + y^2}{x^2 + y^2} &= \lim_{(x, y) \to (0, 0)} \left(1 + \frac{1}{x^2 + y^2}\right) \\
&\overset{r^2 = x^2 + y^2}{=} \lim_{r \to 0} \left(1 + \frac{1}{r^2}\right) \\
&= +\infty.
\end{align*}

(3)
\begin{align*}
\lim_{(x, y) \to (0, 0)} \frac{\sqrt{1 + xy} - 1}{xy} &= \lim_{(x, y) \to (0, 0)} \frac{(\sqrt{1 + xy} - 1)(\sqrt{1 + xy} + 1)}{xy(\sqrt{1 + xy} + 1)} \\
&= \lim_{(x, y) \to (0, 0)} \frac{1}{\sqrt{1 + xy} + 1} \\
&= \frac{1}{2}.
\end{align*}

(6)
\begin{align*}
\lim_{(x, y) \to (0, 0)} \frac{\sin(x^3 + y^3)}{x^2 + y^2} &= \lim_{(x, y) \to (0, 0)} \frac{\sin(x^3 + y^3)}{x^3 + y^3} \cdot \frac{x^3 + y^3}{x^2 + y^2} \\
&= \lim_{(x, y) \to (0, 0)} \frac{\sin(x^3 + y^3)}{x^3 + y^3} \cdot \lim_{(x, y) \to (0, 0)} \frac{x^3 + y^3}{x^2 + y^2} \\
&= 0.
\end{align*}

(8)
\begin{align*}
0 \leq \lim_{\substack{x \to +\infty \\ y \to +\infty}} (x^2 + y^2)\e^{-(x + y)} &\leq \lim_{\substack{x \to +\infty \\ y \to +\infty}} (x^2 + y^2) \lim_{\substack{x \to +\infty \\ y \to +\infty}} \e^{-\sqrt{x^2 + y^2}} \\
&\overset{\substack{x = r\cos\theta \\ y = r\sin\theta}}{=} \lim_{r \to +\infty} r^2 \e^{-r} = 0.\\
\end{align*}
故\(\lim\limits_{\substack{x \to +\infty \\ y \to +\infty}} (x^2 + y^2)\e^{-(x + y)} = 0.\)

\textbf{8.}讨论下列函数在原点的二重极限和二次极限：

(2)\(f(x, y) = \dfrac{x^2(1 + x^2) - y^2(1 + y^2)}{x^2 + y^2}\).

解：
\[\lim_{\substack{x \to 0 \\ y = kx}}f(x, y) = \lim_{x \to 0}\dfrac{x^2(1 + x^2) - k^2x^2(1 + k^2x^2)}{x^2 + k^2x^2} = \dfrac{1 - k^2}{1 + k^2},\]
故二重极限不存在.
\[\lim_{y \to 0} \lim_{x \to 0} f(x, y) = \lim_{y \to 0} \dfrac{0 - y^2(1 + y^2)}{0 + y^2} = -1.\]
\[\lim_{x \to 0} \lim_{y \to 0} f(x, y) = \lim_{x \to 0} \dfrac{x^2(1 + x^2) - 0}{x^2 + 0} = 1.\]

(3)\(f(x, y) = x\sin\dfrac{1}{y} + y\sin\dfrac{1}{x}\).

解：
\[\lim_{\substack{x \to 0 \\ y \to 0}} f(x, y) = \lim_{\substack{x \to 0 \\ y \to 0}} \left(x\sin\dfrac{1}{y} + y\sin\dfrac{1}{x}\right) \leq \lim\limits_{\substack{x \to 0 \\ y \to 0}} |x| + |y| = 0.\]
由\(\lim\limits_{x \to 0} y\sin\dfrac{1}{x}\)和\(\lim\limits_{y \to 0} x\sin\dfrac{1}{y}\)均不存在，故二次极限不存在.

\textbf{9.}验证函数
\[f(x, y) = \begin{cases}
\dfrac{2}{x^2}\left(y - \dfrac{1}{2}x^2\right), & x > 0\text{且}\dfrac{1}{2}x^2 < y \leq x^2, \\
\dfrac{1}{x^2}(2x^2 - y), & x > 0\text{且}x^2 < y < 2x^2, \\
0, & \text{其他点}
\end{cases}\]
在原点不连续，而在其他点连续.

解：当\(x > 0\)且\(\dfrac{1}{2}x^2 < y \leq x^2\)时，
\[\lim_{\substack{x \to 0 \\ y = kx^2}} f(x, y) = \lim_{x \to 0} \dfrac{2}{x^2}\left(kx^2 - \dfrac{1}{2}x^2\right) = 2k - 1,\]
故二重极限不存在，所以在原点不连续.

只需考虑在\(y = \dfrac{1}{2}x^2, y = x^2, y = 2x^2\)，其中\(x > 0\)上的连续性.设\(x_0 > 0\)，则
\[\lim_{\substack{(x, y) \to (x_0, \frac{1}{2}x_0^2) \\ y > \frac{1}{2}x^2}} f(x, y) = \lim_{(x, y) \to (x_0, \frac{1}{2}x_0^2)} \dfrac{2}{x^2}\left(y - \dfrac{1}{2}x^2\right) = 0 = f(x_0, \dfrac{1}{2}x_0^2),\]
\[\lim_{\substack{(x, y) \to (x_0, \frac{1}{2}x_0^2) \\ y \leq \frac{1}{2}x^2}} f(x, y) = 0.\]
故\(f(x, y)\)在\(y = \dfrac{1}{2}x^2\)上连续.
\[\lim_{\substack{(x, y) \to (x_0, x_0^2) \\ y \leq x^2}} f(x, y) = \lim_{(x, y) \to (x_0, x_0^2)} \dfrac{2}{x^2}\left(y - \dfrac{1}{2}x^2\right) = 1 = f(x_0, x_0^2),\]
\[\lim_{\substack{(x, y) \to (x_0, x_0^2) \\ x^2 < y < 2x^2}} f(x, y) = \lim_{(x, y) \to (x_0, x_0^2)} \dfrac{1}{x^2}(2x^2 - y) = 1.\]
故\(f(x, y)\)在\(y = x^2\)上连续.
\[\lim_{\substack{(x, y) \to (x_0, 2x_0^2) \\ x^2 < y < 2x^2}} f(x, y) = \lim_{(x, y) \to (x_0, 2x_0^2)} \dfrac{1}{x^2}(2x^2 - y) = 0 = f(x_0, 2x_0^2),\]
\[\lim_{\substack{(x, y) \to (x_0, 2x_0^2) \\ y \geq 2x^2}} f(x, y) = 0.\]
故\(f(x, y)\)在\(y = 2x^2\)上连续.

综上，\(f(x, y)\)在除原点外的其他点均连续.

\textbf{11.}设\(f(t)\)在区间\((a, b)\)上具有连续导数，\(D = (a, b) \times (a, b)\).定义\(D\)上的函数
\[F(x, y) = \begin{cases}
\dfrac{f(x) - f(y)}{x - y}, & x \neq y, \\
f'(x), & x = y.
\end{cases}\]
证明：对于任何\(c \in (a, b)\)成立
\[\lim_{(x, y) \to (c, c)} F(x, y) = f'(c),\]

解：由拉格朗日中值定理，\(\exists \xi\)介于\(x, y\)之间使得\(f(x) - f(y) = f'(\xi)(x - y)\).
\[\lim_{\substack{(x, y) \to (c, c) \\ x \neq y}} F(x, y) = \lim_{(x, y) \to (c, c)} \frac{f(x) - f(y)}{x - y} = f'(\xi) = f'(c),\]
\[\lim_{\substack{(x, y) \to (c, c) \\ x = y}} F(x, y) = \lim_{(x, y) \to (c, c)} f'(x) = f'(c).\]
综上，\(\lim\limits_{(x, y) \to (c, c)} F(x, y) = f'(c)\).

\textbf{13.}证明：若\(f\)和\(g\)是\(D \subset \RR^n\)上的连续映射，则映射\(f + g\)与函数\(\langle f, g \rangle\)在\(D\)上都是连续的.

解：\(\forall x \in D, \lim\limits_{x \to x_0} (f + g)(x) = \lim\limits_{x \to x_0} f(x) + \lim\limits_{x \to x_0} g(x) = f(x_0) + g(x_0) = (f + g)(x_0)\).

\(\forall x \in D, \lim\limits_{x \to x_0} \langle f, g \rangle (x) = \langle \lim\limits_{x \to x_0} f(x), \lim\limits_{x \to x_0} g(x) \rangle = \langle f(x_0), g(x_0) \rangle = \langle f, g \rangle (x_0)\).

\section*{11.3}
\textbf{1.}设\(D \subset \RR^n\)，\(f: D \to \RR^m\)为连续映射. 如果\(D\)中的点列\(\{x_k\}\)满足\(\lim\limits_{k \to \infty} x_k = a\)，且\(a \in D\)，证明
\[\lim_{k \to \infty} f(x_k) = f(a).\]

解：\(\forall \varepsilon > 0\)，由\(f\)在\(D\)上连续，\(\exists \delta > 0\)，当\(\|x - a\| < \delta\)时，有\(\|f(x) - f(a)\| < \varepsilon\).又由\(\lim\limits_{k \to \infty} x_k = a\)，\(\exists N > 0\)，当\(k > N\)时，有\(\|x_k - a\| < \delta\)，所以当\(k > N\)时，有
\[\|f(x_k) - f(a)\| < \varepsilon.\]
故\(\lim\limits_{k \to \infty} f(x_k) = f(a)\).

\textbf{2.}设\(f\)是\(\RR^n\)上的连续函数，\(c\)为实数. 设
\[A_c = \{x \in \RR^n | f(x) < c\}, \quad B_c = \{x \in \RR^n | f(x) \leq c\}.\]
证明：\(A_c\)为\(\RR^n\)上的开集，\(B_c\)为\(\RR^n\)上的闭集.

解：设\(x_0 \in A_c\)，则\(f(x_0) < c\)，令\(\varepsilon = c - f(x_0) > 0\)，由\(f\)在\(\RR^n\)上连续，\(\exists \delta > 0\)，当\(\|x - x_0\| < \delta\)时，有\(|f(x) - f(x_0)| < \varepsilon\)，即
\[f(x) < f(x_0) + \varepsilon = c.\]
所以\(x \in A_c\)，故\(A_c\)为\(\RR^n\)上的开集.

设\(x_0\)为\(B_c\)的任意聚点，则\(\forall \varepsilon > 0\)，\(\exists x \in B_c\)，使得\(\|x - x_0\| < \varepsilon\)，即\(f(x) \leq c\).由\(f\)在\(\RR^n\)上连续，得
\[\lim_{x \to x_0} f(x) = f(x_0) \leq c.\]
所以\(x_0 \in B_c\)，故\(B_c\)为\(\RR^n\)上的闭集.

\textbf{3.}设二元函数
\[f(x, y) = \dfrac{1}{1 - xy}, \quad (x, y) \in D = [0, 1) \times [0, 1),\]
证明：\(f\)在\(D\)上连续，但不一致连续.

解：由初等函数的连续性，\(f\)在\(D\)上连续. \(n \to \infty\)时，
\[\left|(1 - \frac{1}{n}, 1 - \frac{1}{n}) - (1 - \frac{1}{2n}, 1 - \frac{1}{2n})\right| \to 0,\]
\[\left|f(1 - \frac{1}{n}, 1 - \frac{1}{n}) - f(1 - \frac{1}{2n}, 1 - \frac{1}{2n})\right| = \left|\frac{(4n - 3)n^2}{(4n - 1)(2n - 1)}\right| \to \infty.\]
故\(f\)在\(D\)上不一致连续.

\textbf{6.}设\(f\)是\(\RR^n\)上的连续函数，满足

(1)当\(x \neq 0\)时成立\(f(x) > 0\);

(2)对于任意\(x\)与\(c > 0\)，成立\(f(cx) = cf(x)\)，

证明：存在\(a > 0, b > 0\)使得
\[a|x| \leq f(x) \leq b|x|.\]

解：由\(f\)在\(\RR^n\)上连续，得\(f\)在单位球面\(S = \{x | |x| = 1\}\)上连续，而\(S\)为闭有界集，故\(f\)在\(S\)上必有界，设\(b = \max\limits_{x \in S} f(x)\)，则对任意\(x \neq 0\)，有
\[f(x) = f\left(\frac{x}{|x|}|x|\right) = |x|f\left(\frac{x}{|x|}\right) \leq b|x|.\]
又由\(f\)在\(S\)上连续，且\(f(x) > 0\)，故\(f\)在\(S\)上必有最小值，设\(a = \min\limits_{x \in S} f(x)\)，则对任意\(x \neq 0\)，有
\[f(x) = f\left(\frac{x}{|x|}|x|\right) = |x|f\left(\frac{x}{|x|}\right) \geq a|x|.\]
综上，存在\(a > 0, b > 0\)使得
\[a|x| \leq f(x) \leq b|x|.\]

\textbf{8.}设\(f\)是有界开区域\(D \subset \RR^2\)上的一致连续函数，证明：

(1)可以将\(f\)连续延拓到\(D\)的边界上，即存在定义在\(\bar{D}\)上的连续函数\(\tilde{f}\)，使得\(\tilde{f}|_D = f\)；

(2)\(f\)在\(D\)上有界.

解：(1)设\(x_0 \in \partial D\)，则\(\forall \varepsilon > 0\)，由\(f\)在\(D\)上一致连续，\(\exists \delta > 0\)，当\(\|x - y\| < \delta\)时，有\(|f(x) - f(y)| < \varepsilon\).又由\(x_0\)为\(D\)的边界点，\(\exists x \in D\)，使得\(\|x - x_0\| < \delta\)，所以对任意\(y \in D\)，当\(\|y - x_0\| < \delta\)时，有
\[\|x - y\| \leq \|x - x_0\| + \|y - x_0\| < 2\delta,\]
\[|f(x) - f(y)| < \varepsilon.\]
所以\(\lim\limits_{y \to x_0} f(y)\)存在，记为\(f(x_0)\)，则\(\tilde{f}(x) = f(x)\)当\(x \in D\)，\(\tilde{f}(x) = f(x_0)\)当\(x = x_0\)，则\(\tilde{f}\)在\(D\)上连续.

(2)由\(f\)在\(D\)上一致连续，\(\exists \delta > 0\)，当\(\|x - y\| < \delta\)时，有\(|f(x) - f(y)| < 1\).又由\(D\)为有界开区域，得\(D \subset O(O, R)\)，其中\(O(O, R) = \{x | |x| < R\}\).取\(x_1, x_2, \cdots, x_m \in D\)，使得
\[D \subset \bigcup_{i = 1}^m O(x_i, \delta),\]
则对任意\(x \in D\)，存在\(i\in \{1, 2, \cdots, m\}\)，使得\(\|x - x_i\| < \delta\)，所以
\[|f(x)| \leq |f(x) - f(x_i)| + |f(x_i)| < 1 + \max_{1 \leq i \leq m} |f(x_i)|.\]
\end{document}